% ! TeX root = ./main.tex

\section{Why Telemetry-driven Placement?}\label{sec:amr:bg}

Block-structured AMR simulations running on O(100K) cores face significant performance challenges
due to computational variability~\cite{VanS2009:AMR,Dubey2014:Survey,Schor2018:AMR}.
Frequent mesh refinement necessitates periodic redistribution of work across compute resources.
Stragglers, if left unaddressed, are particularly expensive due to frequent global synchronization
operations---a single straggler can block the entire simulation at synchronization points,
leading to poor cluster utilization~\cite{Petri2003:Missing}.

% ~\cite{Klenk2017:OverviewMPICharacteristics}.

% Two complementary approaches can improve cluster utilization for workloads with computational
% variability --- balancing work among ranks to reduce variability, and asynchronous execution
% strategies to mask residual variability.
Two complementary approaches exist for computational variability---balancing work among ranks to
reduce variability, and asynchronous execution strategies to mask residual variability.
While asynchronous runtimes~\cite{Acun2014:Charm++, Meng2010:Uintah, Daiss2024:HPX,
  Bauer2012:Legion} are commonly used to mask variability, significant time is lost waiting within
MPI routines despite their widespread use ($>$60\% at 1,000
ranks~\cite{Klenk2017:OverviewMPICharacteristics}).
% (see% \S\ref{sec:amr:bg:amr})
Addressing variability via placement is difficult---frameworks often lack meaningful per-block
cost inputs, as predictive models are hard to create and codes often assume uniform costs.
% Addressing variability via placement remains challenging --- while frameworks allow specifying
% per-block costs, these hooks are rarely used meaningfully as predictive cost models do not exist
% in practice. Codes typically use either simple heuristics or assume uniform block costs, ignoring
% dynamic variations.
This gap motivates leveraging runtime telemetry to directly measure workload behavior.

This section provides background for the rest of this paper, beginning with an overview of AMR
codes and current placement techniques in \S\ref{sec:amr:bg:amr}, followed by computational
characteristics of AMR operations in \S\ref{sec:amr:bg:ops}.

% provide hooks for incorporating variable block costs, leveraging these hooks requires cost models
% to predict block costs. Such models require domain scientists to reason about relationships
% between physical properties and computational work, which may not be obvious across problem
% regimes. \ga{Do such cost models exist currently? If so, say so and cite them. If not, be
% explicit.} Runtime telemetry offers the potential to capture fine-grained runtime variations and
% inform placement decisions without relying on complex predictive models. Our work examines the
% potential for telemetry-informed placement to complement asynchronous execution strategies.

% \subsection{Managing Compute Variability in AMR} \label{sec:amr:bg:moti} %\paragraph{Challenge:
% Handling Compute Variability.} Block-structured AMR simulations running on O(100K) cores face
% significant performance challenges due to computational
% variability~\cite{VanS2009:AMR,Dubey2014:Survey,Schor2018:AMR}. Frequent mesh refinement and
% coarsening operations necessitate periodic redistribution of work across compute resources,
% requiring placement policies that can efficiently handle this variability. Stragglers, if left
% unaddressed, are particularly expensive due to frequent global synchronization operations --- a
% single straggler can block the entire simulation at synchronization points, leading to poor
% cluster utilization~\cite{Klenk2017:OverviewMPICharacteristics}. \paragraph{Approaches to
% Computational Variability} Two complementary approaches are theoretically possible to improve
% cluster utilization for workloads with computational variability --- balancing work among ranks to
% reduce total work variability, and asynchronous execution strategies to mask residual variability
% by using idle time to perform useful work. While considerable prior work exists on using
% asynchronous runtimes to mask variability (see \S\ref{sec:amr:bg:amr}), significant cluster
% underutilization attributed to load imbalance persists despite their widespread
% use~\cite{Klenk2017:OverviewMPICharacteristics}. Addressing variability via placement remains
% challenging --- while many frameworks provide hooks for incorporating variable block costs,
% leveraging these hooks requires cost models to predict block costs. Developing such cost models
% requires domain scientists to reason about the relationship between physical properties and
% computational work --- characteristics that may not be obvious or consistent across problem
% regimes. Our work examines the potential for telemetry-informed placement to complement
% asynchronous execution strategies. \parahead{Need for Runtime Telemetry} This gap motivates our
% investigation of telemetry-driven placement as a tool to augment existing placement
% infrastructure. Rather than %ad-hoc cost heuristics, requiring complex predictive models, this
% approach directly measures runtime behavior to guide placement decisions. However, effectively
% collecting and utilizing such measurements, particularly at scale where system noise and
% variability dominate, requires careful consideration of both measurement methodology and placement
% strategies. The next subsections establish the context needed to understand these challenges.
% \ga{I don't know what this paragraph adds, and you also mention cost models again that you provide
% no citation for. Make sure you don't get too philosophical.} Statically estimating these costs is
% challenging --- it requires domain scientists to reason about the relationship between physical
% properties and computational work, characteristics that may not be obvious or consistent across
% problem regimes. This limits the adoption of hooks to incorporate costs for frameworks that do
% provide such hooks. While frameworks provide hooks for incorporating cost models based on proxy
% metrics~\cite{Burst2011:p4est} like refinement levels or solution gradients, leveraging these
% capabilities is challenging. . The challenge lies in the complexity of the relationship between
% physical properties and computational work - it requires domain scientists to reason about
% performance characteristics that may not be obvious or consistent across problem regimes. This has
% led to a gap between AMR's theoretical benefits and achieved performance in practice. 

\subsection{Block-based AMR: Infrastructure and Execution}\label{sec:amr:bg:amr}

% This section describes block-based AMR codes, a sub-class of structured AMR codes, along with
% their execution model and placement mechanisms. While other structured AMR codes share similar
% challenges, we focus our discussion on block-based AMR for simplicity. 
\parahead{Structured AMR Codes}
Structured AMR implementations organize computational grids hierarchically while maintaining a
logically Cartesian structure.
Implementations typically fall into two categories: \emph{block-based} and \emph{patch-based}.
Block-based AMR partitions the domain into uniform sized blocks at each refinement level, managed
using octree data structures.
This approach enables efficient parallel implementations through simplified mesh management and
communication patterns.
Examples include Parthenon~\cite{Grete2023:Parthenon},
Enzo-E/Cello~\cite{Emeri2019:EnzoECelloProject}, and ALPS~\cite{Burst2008:ALPS}.
Patch-based AMR, in contrast, refines regions using arbitrary rectangular patches rather than fixed
blocks, allowing finer control over refinement but requires more complex load-balancing (e.g.
AMReX~\cite{Zhang2021:AMReX}, SAMRAI~\cite{Gunne2016:SAMRAI}).

% . Frameworks such as AMReX~\cite{Zhang2021:AMReX} and SAMRAI~\cite{Gunne2016:SAMRAI} follow this
% approach. %Block- and patch-based AMR share challenges in workload estimation and communication
% locality. The techniques we explore apply to both, but for simplicity, we focus on block-based AMR
% in this paper.

% generally fall into two categories: block-based approaches using octrees, and patch-based
% approaches. Block-based codes partition the domain into uniform-sized blocks at each refinement
% level, managed through octree data structures. This approach enables efficient parallel
% implementations through simplified mesh management and communication patterns.

% In contrast, patch-based frameworks like AMReX, Cactus, Uintah, and FLASH~\cite{Dubey2014:Survey}
% allow for more flexible refinement regions. While our work focuses on block-based implementations,
% many of our findings apply to both approaches. 
\parahead{Execution Model}
Block-based AMR codes operate on structured meshes that are decomposed into mesh blocks, with
multiple blocks typically assigned to each simulation rank.
As the simulation tracks evolving physical phenomena, blocks may be refined or coarsened to adapt
mesh resolution, requiring periodic redistribution across ranks.
Computationally, operations on each block are structured as a \emph{directed acyclic graph} (DAG)
of tasks---including physics computations, mesh management operations, and inter-block
communication.
Beyond these per-block operations, AMR codes also invoke global operations for redistribution and
synchronization, whose characteristics we detail in~\S\ref{sec:amr:bg:ops}.

AMR codes use task-based runtimes to execute DAGs using asynchronous execution strategies.
Some frameworks like Parthenon implement these abstractions directly, while others use
general-purpose runtimes such as Charm++~\cite{Acun2014:Charm++}, Uintah~\cite{Meng2010:Uintah},
HPX~\cite{Daiss2024:HPX}, and Legion~\cite{Bauer2012:Legion}.
Our work incorporating telemetry in work placement complements execution strategies masking
residual imbalance.

% In the common case, the simulation operates on blocks assigned to each rank. Computationally,
% block operations within each timestep are structured as an identical \emph{directed acyclic graph
% (DAG)} of tasks. Every rank executes this DAG every timestep, for each block assigned to it. These
% tasks include compute kernels for physics, mesh management operations, and communication between
% blocks. In addition to these per-block operations, AMR codes also periodically invoke global
% operations for redistribution and synchronization. We describe their characteristics in detail in
% \S\ref{sec:amr:bg:ops}. As the simulation progresses, the number of blocks may grow or shrink over
% time due to (de)refinement operations, necessitating \emph{redistribution}.

\parahead{Placement Mechanisms}
When a redistribution is invoked, placement mechanisms compute new block-rank assignments, and
blocks are migrated accordingly.
Modern block-based AMR codes rely on octree mesh structures combined with space-filling curves
(SFCs) to enable efficient parallel mesh management while approximately preserving communication
locality.
SFC-based partitioning balances uniform block counts effectively but does not account for
computational variability between blocks---a gap we address in this work.
Octree structures, SFC ordering, and their implications for placement flexibility are discussed
further in \S\ref{sec:amr:design} where we present our placement policies.

% Octree-based mesh structures combined with space-filling curves (SFCs) form the foundation of
% modern block-based AMR codes. This approach enables efficient parallel mesh management while
% preserving communication locality through SFC-based block ordering. SFC-based partitioning
% balances uniform block counts well, but must be augmented with mechanisms to handle computational
% variability between blocks (focus of this work). We discuss octrees and SFCs in detail in
% \S\ref{sec:amr:design}. \ga{Forward-references are not fun. Try to eliminate where possible, by either
% removing or mentioning what is relevant and explaining that more details on that will be provided
% when necessary. } % \red{some redudancy here}. Block-based AMR codes operate on structured meshes
% that are decomposed into \emph{mesh blocks} and assigned to simulation ranks. These blocks are
% hierarchically refined to concentrate computational resources on areas with the highest simulation
% activity. The number of blocks may grow or shrink over time due to (de)refinement operations,
% necessitating \emph{redistribution}. %

% While runtimes can hide some imbalance through asynchronous execution, frequent global
% synchronization operations - potentially multiple per timestep - make these codes highly sensitive
% to stragglers. A single slow rank among 100,000s can block the entire simulation at sync points,
% limiting the achievable computation-communication overlap. The lessons from our telemetry-based
% study are applicable regardless of the runtime choice.

% \subsection{Placement in AMR Codes.} \label{sec:amr:bg:octrees}

% --- arising from physics complexity, adaptive solver iterations, and mesh refinement patterns.
% Despite AMR frameworks acknowledging this variability, systematic approaches to quantifying and
% utilizing these costs remain unexplored. Moreover, the relative impact of different placement
% constraints (like communication locality) on overall application performance is poorly understood.
% As scientific simulations grow in complexity and scale, the need for placement strategies that
% automatically adapt to dynamic behaviors observed at runtime becomes increasingly critical. 
\subsection{Characteristics of Key AMR Operations}
\label{sec:amr:bg:ops}

% Collected comments: \ga{Not a fan of calling this useful work, because it implies all our
% improvements focus on useless work -- try something more inclusive, like 'scientific/compute work'
% (as opposed to 'communication work' and 'synchronization work').}, \ga{Do you mean 'should
% dominate'? Because so far you have been arguing there are evidence that significant amount (60\%)
% of runtime is spent not doing this type of work} \ga{Are we calling it 'compute work' or 'useful
% work'? Use the terms you define consistently throughout.} re: sync, @\ga{spell out here and
% elsewhere}

\parahead{Compute Tasks}
Compute tasks perform the core physical and mesh operations of the simulation.
While operations like refinement tagging often have predictable costs, physics kernels are a major
source of variability.
For instance, regions with steep gradients may require more solver iterations, increasing compute
cost.
Importantly, this cost is not directly correlated to spatial area, as each block contains the same
number of computational points (\emph{cells}) regardless of refinement level.
Ultimately, the execution time of these compute tasks represents the core computational cost to
advance the simulation state, establishing the fundamental lower bound on runtime.

\parahead{Communication Tasks}
Communication tasks create complex inter-block dependencies through \emph{boundary exchanges},
where each block communicates with up to 26 neighbors in 3D (faces, edges, and vertices).
Communication volume depends primarily on the number of physical variables exchanged and neighbor
relationships instead of the refinement level of the blocks.
These exchanges typically use non-blocking MPI operations and are latency-sensitive due to small
message sizes and the potential for blocking downstream tasks.
\emph{Flux correction} follows a similar small-message, peer-to-peer pattern for ensuring consistency of
conserved quantities.

\parahead{Redistribution}
Redistribution operations rebalance blocks across ranks as the simulation evolves, adapting to
changes from mesh refinement or coarsening.
They are triggered when blocks are tagged for refinement based on physical criteria, like solution
gradients exceeding a threshold.
Each redistribution computes new block-to-rank mappings and migrates data.
Because redistribution may be invoked frequently and must complete within tight performance
budgets, placement policies must be fast and capable of handling complex scheduling constraints.
The frequency depends on the underlying physics---some problems require frequent adaptation,
others are more stable.

\parahead{Synchronization}
Synchronization operations can be explicit (MPI barriers) or implicit (blocking collectives).
They inherently expose performance variability by forcing all ranks to wait until the last rank
reaches the synchronization point.
This waiting time often increases with scale as the likelihood of encountering a straggler grows.
While asynchronous collectives are suggested as an alternative~\cite{Hoefl2007:CaseAsyncColl},
synchronous operations remain necessary for certain correctness guarantees, like ensuring that 
all ranks execute the same timestep.